\documentclass[10pt,a4paper]{article}

\usepackage[margin=0.6in]{geometry}
\usepackage{enumitem}
\usepackage{titlesec}
\usepackage{hyperref}
\usepackage{xcolor}
\usepackage[T1]{fontenc}
\usepackage{lmodern}
\usepackage{tabularx}

\definecolor{accent}{HTML}{1A5276}
\hypersetup{colorlinks=true,urlcolor=accent,linkcolor=accent}
\titleformat{\section}{\large\bfseries\color{accent}}{}{0em}{}[\titlerule]
\titlespacing*{\section}{0pt}{8pt}{4pt}
\setlist[itemize]{nosep,left=0pt,label=\textbullet,topsep=2pt}
\pagestyle{empty}
\setlength{\parindent}{0pt}

\begin{document}

\begin{center}
  {\LARGE\bfseries Foad S. Farimani, PhD} \\[3pt]
  {\small Robotics Hardware Engineer\enspace$\mid$\enspace Novel Actuators\enspace$\cdot$\enspace Precision Mechatronics\enspace$\cdot$\enspace Design-for-Manufacture}
  \\[3pt]
  {\footnotesize
    Apeldoorn, Netherlands\enspace$\mid$\enspace
    Dutch citizen\enspace$\mid$\enspace
    +31\,614685010\enspace$\mid$\enspace
    \href{mailto:f.s.farimani@gmail.com}{f.s.farimani@gmail.com}\enspace$\mid$\enspace
    \href{https://linkedin.com/in/fsfarimani}{LinkedIn}\enspace$\mid$\enspace
    \href{https://github.com/Foadsf}{GitHub}
  }
\end{center}

\section{Profile}
I invent robot hardware and make it real. For my PhD I designed a new class of pneumatic stepper actuator: fully 3D-printed, electromagnetically stealth, 850 RPM, with no metal, no electronics, and no lubrication, and safe to run inside an MRI scanner. It earned two patents, IEEE ICRA and RA-L papers, and a Hackaday writeup. Since 2020 I have worked at the sharp end of precision manufacturing for ASML's machines (VDL ETG, then NTS), building the tooling and hardware-in-the-loop rigs that turn the hardest mechanical designs into buildable, validated hardware. I am hands-on across the whole stack: mechanism, actuation, sensing, control, test, and manufacture.

\section{Signature Work}
\begin{itemize}
  \item \textbf{Invented PneuAct}, a new class of MRI-compatible pneumatic stepper actuator. Fully 3D-printed, electromagnetically stealth, 850 RPM, no bearings or seals. Two patents, IEEE ICRA 2018 and RA-L 2020, featured on Hackaday and U-Today.
  \item \textbf{Doubled the fatigue life} of Inconel 718 flexures by tracing a recurring failure to its root cause (EDM recast layer) and fixing it with the right surface treatment. Credited by ASML at their 2024 Technology Conference.
  \item \textbf{Led the team} building ASML's EXE5000 cable-slab integration tooling, from fixtures and molds through supplier and manufacturability closure.
\end{itemize}

\section{Experience}

\textbf{System Lead Engineer} \hfill NTS D\&E, Hengelo \\
{\small\textit{Tier-1 ASML supplier, high-precision manufacturing systems}} \hfill {\small Mar 2022--Present}
\begin{itemize}
  \item Doubled the fatigue life of Inconel 718 flexures by chasing a recurring failure down to its real cause (EDM recast layer) and killing it with the right surface treatment; ASML credited the result at their 2024 Technology Conference.
  \item Built the hardware-in-the-loop rig and mockups that validate safety-critical motion-control logic without ever touching production hardware.
  \item Own the qualification tooling that decides whether ASML modules pass or fail: automated pressure-decay, particle, electrical, and mechanical checks against hard acceptance criteria.
  \item Ran system engineering for the EXE:5200 sensor assembly from design kickoff through release.
  \item Chair the Surface Treatment team, mapping suppliers and processes for electropolishing, plasma polishing, passivation, and precision cleaning.
\end{itemize}

\vspace{4pt}
\textbf{Lead Mechanical Engineer} \hfill VDL ETG T\&D \\
{\small\textit{Precision mechanics, tooling, and ASML equipment development}} \hfill {\small Jul 2020--Feb 2022}
\begin{itemize}
  \item Led a multi-engineer team building ASML's EXE5000 cable-slab integration tooling: fixtures, preforming, molds, and supplier and manufacturability closure.
  \item Designed the micron-level leaf-spring mechanics for YieldStar 500 optical fiducials, where stiffness and geometry set the alignment.
  \item Worked early-stage bellows and force-decoupler concepts (ATCM/VES): pressure-drop, stiffness, and supplier-data trade-offs.
\end{itemize}

\vspace{4pt}
\textbf{PhD Researcher, Surgical Robotics} \hfill University of Twente \\
{\small\textit{MRI-compatible surgical robotics and pneumatic actuation}} \hfill {\small 2016--2020}
\begin{itemize}
  \item Invented and built a new class of MRI-compatible pneumatic stepper actuator: fully 3D-printed, electromagnetically stealth, 850 RPM, no bearings or seals.
  \item Designed and prototyped patient-mounted surgical robots with their sensing and control, validated on custom test benches.
  \item Contributed to the EU Horizon 2020 MURAB consortium on robotic MRI and ultrasound biopsy, with a EUR 4.34M budget.
\end{itemize}

\section{Patents}
\begin{itemize}
  \item "MRI Compatible Surgery Assistant Robot", IR 100300442.
  \item "Drugs Injection Smart Robot", IR 100300872.
\end{itemize}

\section{Selected Publications}
\begin{itemize}
  \item "PneuAct-II: Hybrid Manufactured Electromagnetically Stealth Pneumatic Stepper Actuator", IEEE Robotics and Automation Letters, 2020. doi:10.1109/LRA.2020.2974714
  \item "Introducing PneuAct: Parametrically-Designed MRI-Compatible Pneumatic Stepper Actuator", IEEE ICRA, 2018. doi:10.1109/ICRA.2018.8461116
  \item "Frictional characteristics of Fusion Deposition Modeling (FDM) manufactured surfaces", Rapid Prototyping Journal, 2020. doi:10.1108/RPJ-06-2019-0171
  \item "Magnetic Resonance Compatible Pneumatic Actuators for Surgical Robotics", PhD thesis, University of Twente.
\end{itemize}

\section{Skills}
{\small
\begin{tabularx}{\textwidth}{@{}lX@{}}
  \textbf{Robotics \& actuation:} & Novel actuator and mechanism design, pneumatic stepper actuators, compliant mechanisms, patient-mounted robots, sensing, motion control \\
  \textbf{Hardware \& manufacturing:} & Precision mechanics, design-for-manufacture, fixtures, molds, integration tooling, supplier and manufacturability closure, rapid prototyping, 3D printing \\
  \textbf{Test \& validation:} & Hardware-in-the-loop, qualification tooling, pressure decay, particle and acceptance testing \\
  \textbf{Software:} & Python, C\#/.NET, MATLAB, Git, data analysis \\
  \textbf{CAD/CAE:} & Siemens NX, SolidWorks, CATIA, ANSYS/APDL \\
\end{tabularx}
}

\section{Education}
\textbf{PhD, Surgical Robotics}, University of Twente, Enschede \hfill 2016--2020 \\[2pt]
\textbf{MSc, Biomedical Engineering}, Amirkabir University of Technology, Tehran \hfill 2006--2009 \\[2pt]
\textbf{BSc, Mechanical Engineering}, Sharif University of Technology, Tehran \hfill 2000--2006

\end{document}
